\documentclass[11pt,a4paper]{article}
\usepackage[utf8]{inputenc}
\usepackage[T1]{fontenc}
\usepackage{amsmath,amssymb}
\usepackage{geometry}
\usepackage{enumitem}
\usepackage{booktabs}

\geometry{margin=1in}

\title{Interpreting Weighted Sums in a Histogram}
\author{}
\date{}

\begin{document}

\maketitle

Summing $(\text{count in bin} \times \text{bin value})$ gives you the \textbf{approximated total sum of all individual data points} in your dataset.

If you divide that total by the sum of all counts (the total number of pixels or observations), you get the \textbf{binned weighted mean (average)} of your dataset.

\section*{1. What the Formulas Represent}

\begin{itemize}[leftmargin=1.5em]
    \item \textbf{Sum of Products:}
    \begin{equation}
        \sum_{i=1}^{k} \left(\text{Count}_i \times \text{Bin Value}_i\right) \approx \text{Total Sum of Raw Values}
    \end{equation}
    \textit{In NDVI terms:} If you have $1{,}000$ pixels and each pixel represents a $10\text{ m} \times 10\text{ m}$ ($100\text{ m}^2$) area, multiplying that sum by the pixel area gives you the \textbf{integrated NDVI} (a proxy for total landscape greenness or total biomass across the entire region).

    \item \textbf{Divided by Total Pixels:}
    \begin{equation}
        \frac{\sum_{i=1}^{k} \left(\text{Count}_i \times \text{Bin Value}_i\right)}{\sum_{i=1}^{k} \text{Count}_i} \approx \text{Mean NDVI}
    \end{equation}
    \textit{In NDVI terms:} This gives you the \textbf{average NDVI} of your entire image or field from the histogram alone, without needing the original raw raster pixels.
\end{itemize}

\section*{2. Why ``Bin Value'' Choice Matters}

Because a bin covers a continuous range (e.g., $[0.40, 0.50]$), multiplying by different parts of the bin changes your accuracy:

\begin{itemize}[leftmargin=1.5em]
    \item \textbf{Bin Midpoint (Recommended):} Using the center of the bin ($0.45$) assumes values are evenly spread inside the bin. Errors tend to cancel out, giving the closest estimate to the true raw sum and mean.
    \item \textbf{Lower Bin Edge ($0.40$):} Systematically \textbf{underestimates} the true sum and mean.
    \item \textbf{Upper Bin Edge ($0.50$):} Systematically \textbf{overestimates} the true sum and mean.
\end{itemize}

\section*{3. Quick Numerical Example}

Suppose a tiny $10$-pixel field has the following binned counts:

\begin{center}
\begin{tabular}{@{}cccc@{}}
\toprule
\textbf{NDVI Bin} & \textbf{Bin Midpoint} & \textbf{Pixel Count} & \textbf{Count $\times$ Midpoint} \\
\midrule
$[0.1, 0.3]$ & $0.2$ & $2$ & $2 \times 0.2 = 0.4$ \\
$[0.5, 0.7]$ & $0.6$ & $8$ & $8 \times 0.6 = 4.8$ \\
\midrule
\textbf{Total} & --- & $\mathbf{10}$ & $\mathbf{5.2}$ \\
\bottomrule
\end{tabular}
\end{center}

\begin{enumerate}[leftmargin=1.5em]
    \item \textbf{Sum of $(\text{count} \times \text{midpoint})$:} $0.4 + 4.8 = \mathbf{5.2}$ (Approximated total NDVI sum).
    \item \textbf{Divide by total count ($10$):} $5.2 / 10 = \mathbf{0.52}$ (Approximated average NDVI of the field).
\end{enumerate}

\end{document}
